\chapter{Datasets and Experimental Setup}\label{ch:setup}

\section{Hardware and software}

All experiments ran on one workstation (\cref{tab:env}). The binding resource was the CPU, not the GPU. One Spark-Scala
program takes 40--46~s in the sandbox, so Algorithm~3 for one round took about 2.0 to 5.1 hours at 36 parallel workers
(\cref{tab:gen-timings}). The longer values include waiting for the shared validation lock when several seeds ran at
once. Training and serving fit comfortably on one 24~GB GPU each.

{\small
\rowcolors{2}{rowshade}{white}
\begin{longtable}{@{}p{3.6cm}p{11.0cm}@{}}
\caption{Computing environment. Versions are those installed in the project's virtual environment at the time of
writing; container images are pinned by digest.}\label{tab:env}\\
\toprule\rowcolor{headshade}Component & Value\\\midrule\endfirsthead
\toprule\rowcolor{headshade}Component & Value\\\midrule\endhead
\bottomrule\endfoot
CPU & Intel Xeon w7-3455, 40 logical CPUs\\
Memory & 102 GB\\
GPUs & 2 $\times$ NVIDIA RTX 4500 Ada Generation, 24,570 MiB each (driver 595.97). GPU~1 also drives the display and
has about 1.75~GB less free memory\\
Operating system & Ubuntu 24.04 under WSL2 on Windows; all data on the Linux file system\\
Python stack & Python 3.12.3; PyTorch 2.14.0; vLLM 0.33.0 (serving); Transformers 5.20.2; TRL 1.17.1; PEFT 0.21.2;
pandas 3.1.0; NumPy 2.3.5; sqlglot 30.29.0; MLflow-skinny 3.18.0 (local tracking, telemetry off)\\
Code sandboxes & Docker 29.9.3. Spark: \texttt{apache/spark:3.5.7} by digest, open-file limit raised to 4,096.
Python: \texttt{python:3.11-slim} by digest. R: \texttt{r-base:4.4.2} by digest plus Debian's \texttt{r-cran-dplyr}
(R 4.6.1, dplyr 1.2.1 after installation). No network, read-only file system, non-root user\\
SQL & SQLite on the host, read-only connection, SELECT only, 5~s limit per query\\
Report & TeX Live 2025 in Docker (\texttt{texlive/texlive:TL2025-historic}, pinned by digest)\\
\end{longtable}
}

\section{Models}

{\small
\rowcolors{2}{rowshade}{white}
\begin{longtable}{@{}p{3.4cm}p{2.1cm}p{9.1cm}@{}}
\caption{Models used. Revisions are pinned Hugging Face commits; all four weight-file checksums were verified after
download.}\label{tab:models}\\
\toprule\rowcolor{headshade}Model & Revision & Role\\\midrule\endfirsthead
\toprule\rowcolor{headshade}Model & Revision & Role\\\midrule\endhead
\bottomrule\endfoot
Llama-3.1-8B (base)~\cite{grattafiori2024llama3} & \texttt{d04e592} & Base track: questions, SQL and Scala generation,
training and evaluation, through raw completions\\
Llama-3.1-8B-Instruct & \texttt{0e9e39f} & Instruct track: all of the above through the repository's chat call.
Also the repairer in every repair-TS run, including the base track\\
Qwen2.5-7B-Instruct & \texttt{a09a354} & Pipeline testing only; never reported as a result\\
\end{longtable}
}

\section{Training-table corpus}\label{sec:corpus}

\begin{paperbox}[title={Paper}]
``A large corpus'' of real tables, unreleased. The repository's file names suggest web tables, in 3,000-table slices
per round.
\end{paperbox}
\begin{oursbox}
We built the corpus from the GitTables 1M licensed subset~\cite{hulsebos2021gittables,gittableszenodo}
(\fpath{ts/corpus/build.py}; \cref{tab:gen-corpus}). We chose it because it is licensed (CC0 per its Zenodo record),
large, and from a different source (GitHub CSV files) than the Wikipedia-based test sets.
\end{oursbox}

\begin{itemize}
\item \textbf{Shape filters.} We keep tables with 5--30 columns (the repository requires at least 5), 10--200 rows, and
at most 2,000 serialised tokens.
\item \textbf{Quality filters.} We drop tables with blank or duplicate headers, or with more than 50\% empty cells.
\item \textbf{Near-duplicates.} Tables whose header sets have Jaccard similarity $\ge 0.9$ are near-duplicates, and at
most 5 per group are kept. Headers that differ only by case count as duplicates: SQLite ignores case in column names,
and such a table crashed a real run.
\item \textbf{Contamination.} The header set of every kept table is compared with all 1,141 header sets of the four
test sets at Jaccard 0.9. This removed 38 tables.
\end{itemize}

The limit of 5 per group was the smallest that yields the 9,000 tables each seed needs (3 rounds $\times$ 3,000). It
was chosen on table counts alone, before any fine-tuning result existed. Limits of 1, 2, 3, 5 and 10 kept 5,036, 7,430,
8,751, 10,378 and 12,640 tables. The case-insensitive header rule, added later, reduced 10,378 to the final 10,053.

\begin{observation}
The corpus is less Wikipedia-like than the paper's (likely) web tables. Its largest topic accounts for 46.5\% of
tables (\texttt{id\_tables}), and the contamination check compares headers only, so it cannot catch identical data
under different headers.
\end{observation}

{\small
\rowcolors{3}{rowshade}{white}
\setlength{\tabcolsep}{4.2pt}
\begin{longtable}{@{}lp{9.5cm}@{}}
\caption{Training-table corpus built from GitTables (\texttt{ts/corpus/build.py}).}\label{tab:gen-corpus}\\
\toprule
\rowcolor{headshade}Property & Value\\
\midrule
\endfirsthead
\multicolumn{2}{l}{\textit{\tablename~\thetable{} (continued)}}\\
\toprule
\rowcolor{headshade}Property & Value\\
\midrule
\endhead
\midrule
\multicolumn{2}{r}{\textit{continued on next page}}\\
\endfoot
\bottomrule
\endlastfoot
Source & GitTables 1M licensed subset, version 0.0.6, DOI \texttt{10.5281/zenodo.6517052}, licence cc-zero\\
Files / bytes & 96 zips / 16,253,928,856\\
Tables examined & 1,018,649\\
Tables kept & 10,053 (target 20,000; sampling seed 42)\\
Columns & 5--30\\
Rows & 10--200\\
Serialised size & at most 2,000 tokens\\
De-duplication & header Jaccard $\geq$ 0.9, at most 5 tables per header group\\
Contamination check & 1,141 test-set header sets\\
Topics used / largest topic share & 68 / 46.5\%\\
Rejected: too many tokens & 447,282\\
Rejected: columns & 272,371\\
Rejected: rows & 147,845\\
Rejected: duplicate of kept & 78,872\\
Rejected: blank header & 36,823\\
Rejected: mostly empty & 21,402\\
Rejected: duplicate header & 3,958\\
Rejected: matches benchmark & 38\\
Rejected: unreadable & 5\\
\end{longtable}
}
\vspace{-0.6\baselineskip}
{\footnotesize\color{slate}\textit{Note.} Source: \fpath{results/corpus_manifest.json}. Rejection counts are first-failing filter, so they do not sum to the examined total minus the kept total exactly when a table is counted under its first failing rule.\par}
\medskip


\section{Test sets}\label{sec:testsets}

{\small
\rowcolors{2}{rowshade}{white}
\begin{longtable}{@{}p{1.9cm}p{4.0cm}p{1.5cm}p{1.4cm}p{5.5cm}@{}}
\caption{Test sets: how we rebuilt each one, and how it compares with the paper's. Every gold answer was re-checked
against the dataset's own engine.}\label{tab:testsets}\\
\toprule\rowcolor{headshade}Test set & Our construction & Ours & Paper & Match and checks\\\midrule\endfirsthead
\toprule\rowcolor{headshade}Test set & Our construction & Ours & Paper & Match and checks\\\midrule\endhead
\bottomrule\endfoot
WikiSQL~\cite{zhong2017seq2sql,wikisqlrepo} & Test split; first 1,000 of a seed-42 shuffle that fit the prompt limit
(16 skipped). Cells lowercased in prompt and database & 1,000 & 1,000 & Count only; cases probably differ. 0 gold
mismatches against WikiSQL's engine\\
WikiTQ~\cite{pasupat2015wikitq,wikitqrepo} & \texttt{pristine-unseen-tables} (the dataset's test split); same sampling
(64 skipped) & 1,000 & 1,000 & Count only. Gold is an answer list\\
Spider~\cite{yu2018spider} & Test split; gold SQL touching one table (1,230), minus 32 over the prompt limit & 1,198 &
1,198 & \textbf{Count matches exactly}; cases unconfirmed. 0 mismatches, 6 differ only in row order\\
BIRD~\cite{li2023bird} & Dev split (its test set is hidden); gold SQL touching one table. Prompt shows the first 20,
10, 5 or 3 rows; SQL runs on the whole table; evidence strings not shown & 363 & 356 & \textbf{7 more than the paper,
cause unknown}. No single filter gives 7 (empty gold 0, NULL gold 1, failing gold 0, empty tables 0, JOIN in gold 1,
duplicate gold 3). Cases are not deleted to force the count\\
Text2Analysis & Not available (``after the company's approval process'') & -- & 271 & Not possible; excluded from all
averages on both sides\\
AutoTest ST / RT~\cite{autotestrepo} & Columns whose repository prompt fits 70\% of 4,096 tokens; one ST row whose
value list cannot be parsed is dropped & 1,126 / 1,081 & 1,126 / 1,081 & \textbf{Counts match exactly}; 42 and 36
columns contain errors\\
\end{longtable}
}

\paragraph{Dev sets.} Hyper-parameter choices for DPO and self-consistency were made on dev sets only
(\fpath{ts/bench/dev.py}). These are WikiSQL dev (300), Spider dev single-table (575, of which 411 fit the prompt limit)
and WikiTQ \texttt{random-split-1-dev} (300). None overlaps a test set. BIRD has no unused split, so it is not tuned on.

\paragraph{Prompt-table conventions.}
\begin{itemize}
\item Each table is renamed \texttt{"table"} and gold SQL is rewritten to match, as the repository's prompt expects.
\item WikiSQL cells are lowercased in both prompt and database. WikiSQL stores them lowercase, and the repository
lowercases gold SQL. Showing original case instead lowered scores sharply (\cref{sec:fail-gate}).
\end{itemize}

{\small
\rowcolors{3}{rowshade}{white}
\setlength{\tabcolsep}{4.2pt}
\begin{longtable}{@{}llr@{}}
\caption{Saved case and table ID lists, so the exact sets can be rebuilt.}\label{tab:gen-idlists}\\
\toprule
\rowcolor{headshade}Set & ID list & IDs\\
\midrule
\endfirsthead
\multicolumn{3}{l}{\textit{\tablename~\thetable{} (continued)}}\\
\toprule
\rowcolor{headshade}Set & ID list & IDs\\
\midrule
\endhead
\midrule
\multicolumn{3}{r}{\textit{continued on next page}}\\
\endfoot
\bottomrule
\endlastfoot
WikiSQL & \texttt{results/ids/wikisql.txt} & 1000\\
WikiTQ & \texttt{results/ids/wikiqa.txt} & 1000\\
Spider & \texttt{results/ids/spider.txt} & 1230\\
BIRD & \texttt{results/ids/bird.txt} & 363\\
Training corpus & \texttt{results/ids/corpus.txt} & 10053\\
\end{longtable}
}
\vspace{-0.6\baselineskip}
{\footnotesize\color{slate}\textit{Note.} Spider's and WikiTQ's lists may include cases later dropped by the prompt-length limit; the evaluated $n$ is in Table~\ref{tab:gen-vanilla-sql}.\par}
\medskip


\section{Hyper-parameters and settings}\label{sec:hparams}

\Cref{tab:hparams} lists every setting of the core study with its source. The complete configuration file is
reproduced in \cref{app:config}.

{\small
\rowcolors{2}{rowshade}{white}
\begin{longtable}{@{}p{4.3cm}p{5.5cm}p{1.4cm}p{3.4cm}@{}}
\caption{Settings of the NL-to-SQL study and their source: \emph{paper}, \emph{repo} (repository) or \emph{ours} (the
paper is silent).}\label{tab:hparams}\\
\toprule\rowcolor{headshade}Setting & Value & Source & Note\\\midrule\endfirsthead
\toprule\rowcolor{headshade}Setting & Value & Source & Note\\\midrule\endhead
\bottomrule\endfoot
Arms & vanilla, no-validation, Table-Specialist & paper & Section 6.1\\
Seeds & 1, 2, 3 & ours & \\
Rounds & 3 & paper & ``up to 3''\\
Tables per round & 3,000, new each round & paper / repo & ``e.g.\ 3000''; repo 1st/2nd/3rd 3K files\\
Restart each round & from vanilla & paper (partly) & Sec.~6.5: restarting better; main-results choice unstated\\
Training data per round & that round's validated pairs only & paper & Algorithm 1\\
Generator per round & previous round's adapters (vanilla in round 1) & paper & Algorithm 1\\
Second-language model & fine-tuned in lockstep & paper allows & Sec.~5: vanilla or lockstep\\
Question model & vanilla, every round & ours & unstated\\
Question temperature & 0.7 & ours & \\
Code temperature & 0 (greedy) & repo default & \\
Max tokens, base & questions 128, SQL 128, Scala 256 & ours & \\
Max tokens, Instruct & questions 256, SQL 256, Scala 512 & ours & chat replies add prose\\
Request API, base & raw completions, no system message & ours & base has no chat template\\
Request API, Instruct & chat with the repo system message & repo & \fpath{queryWithDeployment.py}\\
Validation samples & 10 at 50\%, random states 0--9 & repo / paper & \fpath{prepare_exec.py}; Algorithm 3\\
Agreement & all 10 samples & paper & Algorithm 3\\
Numeric tolerance (validation) & $10^{-6}$ & ours & \\
Fine-tuning method & LoRA & paper & \\
LoRA rank / alpha / dropout & 16 / 32 / 0.05 & ours & \\
Target modules & all linear layers & ours & \\
Learning rate, schedule & $2\times10^{-4}$, cosine & ours & paper gives only an OpenAI multiplier\\
Epochs & 3 & ours & \\
Batch & 1\% of the examples (gradient accumulation) & paper & \\
Max sequence length & 4,096 tokens & repo & \\
Precision & bf16 & ours & \\
Loss & completion-only & ours & prompts cost no gradient\\
Validation workers & 36 Spark sandboxes, 1 CPU each & ours & one validation at a time machine-wide\\
Evaluation prompts & zero-shot (repo default) and few-shot & repo / ours & \\
Evaluation max tokens & base 128; Instruct 256 & ours & \\
\end{longtable}
}

\section{Random seeds}\label{sec:seeds}

{\small
\rowcolors{2}{rowshade}{white}
\begin{longtable}{@{}p{4.6cm}p{10.0cm}@{}}
\caption{Where randomness enters, and how each source is seeded. ``$s$'' is the study seed (1, 2 or 3); ``$r$'' the
round.}\label{tab:seeds}\\
\toprule\rowcolor{headshade}Source of randomness & Seeding\\\midrule\endfirsthead
\toprule\rowcolor{headshade}Source of randomness & Seeding\\\midrule\endhead
\bottomrule\endfoot
Table order per seed & \texttt{random.Random($s$).shuffle(corpus)}; round $r$ takes positions $3000(r-1)$ to $3000r$\\
Question requirements & repository sampler with RNG $1000s+r$\\
Question sampling (LLM) & request seed $s$, $T=0.7$\\
Code generation & greedy; request seed $s$ (TS loop), 0 (evaluation)\\
Validation row samples & \texttt{DataFrame.sample(frac=0.5, random\_state=$k$)}, $k=0..9$ (fixed across seeds)\\
LoRA training & trainer seed $s$\\
bestof4 / DPO candidates & request seeds $100s+j$, $j=1,2,3$, $T=0.7$\\
Repair retries & request seeds $100s+1..3$, $T=0.7$\\
Same-data control extras & \texttt{random.Random($s$)}; donor seed $s \bmod 3 + 1$\\
Error-detection columns & \texttt{random.Random($1000s+r$)}; errors at $T=0.7$ with request seed $s$\\
Error-detection shuffles & \texttt{random.seed($j$)}, $j=0..4$ (repository)\\
Self-consistency samples & 8 samples at $T=1.0$\\
Bootstrap & 10,000 draws, fixed RNG seed 0\\
Corpus sampling, test-set sampling & seed 42\\
\end{longtable}
}

Vanilla models are run once with greedy decoding. vLLM's batched greedy decoding is not guaranteed to be bit-identical
across runs. Every answer is stored in the response cache, and a rerun reuses it, so the reported results are fixed.

\section{Experiment identifiers}\label{sec:runids}

Every run has an identifier checked against one pattern before any GPU time is used (\fpath{ts/tracking.py}):
\begin{center}
\texttt{(smoke|final|gate|train|baseline)(\_[a-z0-9]+)+}
\end{center}
Training runs are \texttt{train\_<arm>\_<profile>\_seed<s>}, with per-round directories \texttt{round<r>/}. Test cells
are \texttt{<run>\_round<r>\_<testset>\_<prompt>}. Each cell writes \fpath{results/baseline/<id>.json} (summary and
interval) and \fpath{runs/<id>/preds.jsonl} (per-case verdicts). \Cref{app:registry} lists every run family, its
identifiers, seeds and rounds.
